% Einheit 2 von 4 – Äquivalenzumformungen (bank/lineare-gleichungen/e2.jsonl, zusammenbau v0.1)
%% TODO Zweigzeile Teil 1: was hier gelernt wird (Satz in Schülersprache aus dem Einheitstitel) – Einheit 2
\einheitenkopf[e2]{Einheit 2 von 4 · Äquivalenzumformungen}
\zweigzeile{ab Kl. 7, je nach Buch bis Kl. 8 · P10}
\setcounter{aufgabe}{11}

%% TODO Titel: Ich-kann-Satz fehlt in der Bank (Platzhalter: Kettenname) – Nr. 12
%% TODO Anweisung: ein Satz über der ersten Teilaufgabe fehlt in der Bank – Nr. 12
\begin{aufgabe}{Vorzahl und Zeichen bei x lesen}
\begin{teile}
\teil $-5x = 30$ – was steht bei $x$? \leerfeld
\end{teile}
\end{aufgabe}

%% TODO Titel: Ich-kann-Satz fehlt in der Bank (Platzhalter: Kettenname) – Nr. 13
%% TODO Anweisung: ein Satz über der ersten Teilaufgabe fehlt in der Bank – Nr. 13
\begin{aufgabe}{Reihenfolge bestimmen}
\begin{teile}
\teil $3x + 7 = 22$ – welche Umformung zuerst? $\mid$ \leerfeld
\end{teile}
\end{aufgabe}

%% TODO Titel: Ich-kann-Satz fehlt in der Bank (Platzhalter: Kettenname) – Nr. 14
%% TODO Anweisung: ein Satz über der ersten Teilaufgabe fehlt in der Bank – Nr. 14
\begin{aufgabe}{Umformen}
\begin{teile}
\teil $x + 11 = 18$ – was kommt hinter den Strich? $\mid$ \leerfeld
\end{teile}
%% TODO Darstellung für schwach fehlt in der Bank (lineare-gleichungen-e2-k3-s1-v1; Abschnitt „Für schwache Schüler“)
\swz{$x + 8 = 15$}{}
%% TODO Darstellung für schwach fehlt in der Bank (lineare-gleichungen-e2-k3-s1-v2; Abschnitt „Für schwache Schüler“)
\swz{$x - 5 = 12$}{}
%% TODO Darstellung für schwach fehlt in der Bank (lineare-gleichungen-e2-k3-s1-v3; Abschnitt „Für schwache Schüler“)
\swz{$x + 16 = 25$}{}
%% TODO Darstellung für schwach fehlt in der Bank (lineare-gleichungen-e2-k3-s1-v4; Abschnitt „Für schwache Schüler“)
\swz{$x - 11 = 3$}{}
%% TODO Darstellung für schwach fehlt in der Bank (lineare-gleichungen-e2-k3-s1-v5; Abschnitt „Für schwache Schüler“)
\swz{$x - 7 = 30$}{}
\swfrage{Was bleibt gleich, was ändert sich?}
%% TODO Darstellung für schwach fehlt in der Bank (lineare-gleichungen-e2-k3-s2-v1; Abschnitt „Für schwache Schüler“)
\swz{$6x = 54$}{}
%% TODO Darstellung für schwach fehlt in der Bank (lineare-gleichungen-e2-k3-s3-v1; Abschnitt „Für schwache Schüler“)
\swz{$x + 12 = 5$}{}
%% TODO Darstellung für schwach fehlt in der Bank (lineare-gleichungen-e2-k3-s4-v1; Abschnitt „Für schwache Schüler“)
\swz{$\text{$3x + 5 = 26$ – mit Probe}$}{}
%% TODO Darstellung für schwach fehlt in der Bank (lineare-gleichungen-e2-k3-s5-v1; Abschnitt „Für schwache Schüler“)
\swz{$-3x = 18$}{}
%% TODO Darstellung für schwach fehlt in der Bank (lineare-gleichungen-e2-k3-s6-v1; Abschnitt „Für schwache Schüler“)
\swz{$\frac{x}{4} = 6$}{}
%% TODO Darstellung für schwach fehlt in der Bank (lineare-gleichungen-e2-k3-s7-v1; Abschnitt „Für schwache Schüler“)
\swz{Eine zweischrittige Gleichung mit der Lösung 6?}{}
\swz[3]{$\text{$3(x + 7) = 12$ – löse und mache die Probe. \hfill (P10 2020 OS)}$}{}
\end{aufgabe}

%% TODO Titel: Ich-kann-Satz fehlt in der Bank (Platzhalter: Kettenname) – Nr. 15
%% TODO Anweisung: ein Satz über der ersten Teilaufgabe fehlt in der Bank – Nr. 15
\begin{aufgabe}{Umformen – Fehler finden, Begründen}
\swz[3]{Jana rechnet: \rechnung{3x &= 15 &&\mid -3 \\ x &= 12} Finde den Fehler.}{}
\swz[3]{Leo rechnet: \rechnung{x + 8 &= 14 \\ x &= 14 + 8 \\ x &= 22} Finde den Fehler.}{}
\swz[3]{Sara rechnet: \rechnung{3x + 9 &= 21 &&\mid :3 \\ x + 9 &= 7} Finde den Fehler.}{}
\swfrage{Warum darf man bei $x + 4 = 10$ auf beiden Seiten 4 abziehen?}
\end{aufgabe}

\uebersichtskasten{Umformen: Beide Seiten gleich behandeln – dann bleibt die Lösung gleich. \\ +7 ↔ −7 \quad ·4 ↔ :4 \quad x + 7 = 12 | −7 → x = 5 \\ Reihenfolge: erst Plus und Minus wegbringen, dann Mal und Geteilt. \\ 2x + 5 = 13 | −5 → 2x = 8 | :2 → x = 4 \\ Probe: Lösung einsetzen, beide Seiten ausrechnen – gleich? (wA)}
